\section{Cronograma}

\begin{samepage}
O cronograma do trabalho está estruturado em duas etapas principais. A primeira, de natureza exploratória e preparatória, abrange os meses de abril a julho e corresponde à elaboração do projeto de TCC. A segunda etapa, entre agosto e dezembro, contempla o desenvolvimento técnico, a análise experimental, a redação e a apresentação do trabalho final. A Tabela~\ref{tab:cronograma} detalha as atividades planejadas e seus respectivos períodos de execução.

\begin{table}[H]
\centering
\scriptsize
\resizebox{\textwidth}{!}{%
\begin{tabular}{|p{7cm}|c|c|c|c|c|c|c|c|c|}
\hline
\textbf{Atividade} & Abr & Mai & Jun & Jul & Ago & Set & Out & Nov & Dez \\
\hline
Levantamento bibliográfico e estudo teórico & X  & X  &    &    &    &    &    &    &    \\
Definição do escopo e dos objetivos do trabalho & X  & X  &    &    &    &    &    &    &    \\
Redação da introdução, justificativa e objetivos &    & X  & X  &    &    &    &    &    &    \\
Revisão do referencial teórico &    &    & X  & X  &    &    &    &    &    \\
Definição da metodologia e do cronograma do projeto &    &    & X  & X  &    &    &    &    &    \\
Apresentação do projeto à banca &    &    &    & X  &    &    &    &    &    \\
\hline
Implementação do algoritmo genético &    &    &    &    & X  & X  &    &    &    \\
Execução e registro dos experimentos &    &    &    &    &    & X  & X  &    &    \\
Análise dos resultados obtidos &    &    &    &    &    &    & X  & X  &    \\
Redação da versão final do TCC &    &    &    &    &    &    & X  & X  & X  \\
Apresentação final do TCC &    &    &    &    &    &    &    &    & X  \\
\hline
\end{tabular}%
}
\caption{Cronograma de atividades do projeto.}
\label{tab:cronograma}
\end{table}
\end{samepage}
